\documentclass{article}

\usepackage[utf8]{inputenc}
\usepackage[T1]{fontenc}
\usepackage{amsmath,amssymb,amsthm}
\usepackage{mathtools}
\usepackage{hyperref}
\usepackage{booktabs}

\makeatletter
\newcommand{\citep}[2][]{%
  \begingroup
  \def\@tmpkey{#2}%
  \ifx\@tmpkey\@empty\else
    (\ifx\@empty#1\@empty\else #1, \fi
     \@nameuse{bibkey@\@tmpkey})%
  \fi
  \endgroup
}
\expandafter\def\csname bibkey@lasser2026exo\endcsname{Exogenous~Verification}
\expandafter\def\csname bibkey@lasser2026phase\endcsname{Phase~Redundancy}
\expandafter\def\csname bibkey@lasser2026micro\endcsname{Microfoundation}
\expandafter\def\csname bibkey@lasser2026horizon\endcsname{Horizon~Aware}
\makeatother

\newtheorem{definition}{Definition}
\newtheorem{lemma}{Lemma}
\newtheorem{theorem}{Theorem}
\newtheorem{remark}{Remark}
\newtheorem{assumption}{Assumption}

\newcommand{\Aadv}{A_{\mathrm{adv}}}
\newcommand{\deltaadv}{\delta_{\mathrm{adv}}}
\newcommand{\Witness}{\mathcal{W}}
\newcommand{\Ledger}{\mathcal{L}}
\newcommand{\rext}{r_{\mathrm{ext}}}
\newcommand{\mindep}{m_{\mathrm{eff}}^{\mathrm{indep}}}
\newcommand{\deltastar}{\delta_{*}}
\newcommand{\Tbeta}{T_{\beta}}

\newcommand{\Vem}{V_{\mathrm{env}}}
\newcommand{\eenv}{\epsilon_{\mathrm{env}}}
\newcommand{\Bclip}{B}

\title{Lemma 5e (Environment-side witness extension): Full Proof Draft v2}
\author{Paper 10 working draft for codex review}
\date{}

\begin{document}

\maketitle

\section{Goal and changes from v1}

V1 promoted the inline sketch to a structured proof; Round~A
flagged five issues, all addressable.

\textbf{Round A issues + v2 fixes.}
\begin{itemize}
\item \textbf{Q1 (trust model).}  V1's ``$s_{\mathrm{env}} \neq
  s(\mathrm{agent})$'' was insufficient.  V2 explicitly requires
  the witness substrate to be \emph{outside the adversary's
  write-access / failure domain} (not merely a different substrate
  partition), and names the trusted-setup honesty assumption
  (Paper~5's at-least-one-honest participant) as a condition.
\item \textbf{Q2 (KL derivation gaps).}  V2 adds:
  (a) the marginal-to-joint KL step via the data-processing
  inequality (DPI);
  (b) explicit support conventions ($p_0 \in (0,1)$ for Bernoulli;
  $\lambda_0 \geq \eenv > 0$ for Poisson, with non-overlapping
  support handled by (C5.HOEFF) clipping inherited via Lemma~6);
  (c) a two-sided KL floor option for detection classes that
  include any threshold-exceeding shift (not just degradation).
\item \textbf{Q3 ($V_{\mathrm{env}}$ enumeration).}  V2 makes
  $\Vem$ explicitly deployment-policy-derived, with the four
  canonical variables as the illustrative minimum.  Reconciles
  $I_8$ (which lists $\rext, \Delta r_K$, $\lambda$,
  substrate-distinctness) vs.\ deployment-tooling §8.1 (which lists
  $\rext$, distinctness, $\lambda$, trusted-setup): drops
  $\Delta r_K$ from $\Vem$ (it's strategy-dependent, not exogenous;
  $\Delta r_K$ is captured agent-side); adds trusted-setup status
  flag to the canonical environment list.  Integration item to
  update $I_8$'s enumeration accordingly.
\item \textbf{Q4 (R2 connection).}  V2 refines the R2 statement:
  R2 covers manipulations targeting \emph{only} unmonitored
  exogenous variables that produce \emph{no} threshold-exceeding
  shift in any $v \in \Vem$.  Manipulations of unmonitored
  variables that causally shift a monitored $v$ are detected
  through that monitored projection (not R2).
\item \textbf{Q5 (theorem-level conditions).}  V2 introduces
  (C12.ENV-WIT) as a new theorem-level condition (or expanded
  $I_8$ sub-clauses) covering: published fixed $\Vem$,
  per-variable source/witness substrate partition, trust model
  with witness-outside-adversary-domain, setup-honesty assumption,
  per-variable alternative/threshold calibration.  Plus expansion
  of (C5.SPRT) sub-clauses to explicitly cover environment-side
  LLR.
\end{itemize}

\textbf{Carried forward from v1 (unchanged in v2):}
\begin{itemize}
\item Symmetric-construction skeleton (mirror of Paper 5's
  substrate-exclusivity).
\item Per-variable KL floor derivations for the four canonical
  variables (now with v2 support fixes).
\item Combined floor $\deltaadv^{\mathrm{env}} = \min_v
  \delta_v^{\mathrm{env}} > 0$.
\end{itemize}

\section{Setup: environment-side substrate-exclusive witnesses}

\subsection{Symmetric construction (mirror of Paper 5)}

\citep[Exogenous~Verification]{lasser2026exo} constructs
substrate-exclusive witnesses for agent-side observables: a witness
on substrate $s' \neq s(\mathrm{agent})$ attests to events on
substrate $s(\mathrm{agent})$ via cross-substrate verification.
The construction relies on three properties:
\begin{enumerate}
\item \emph{Substrate distinctness:} the witness substrate $s'$ is
  causally separable from the observed substrate $s$.
\item \emph{Cryptographic commitment:} witness records are committed
  via Pedersen commitments \emph{(or analogous binding scheme)}
  ensuring tamper-evidence.
\item \emph{Ledger publication:} commitments are published to the
  verification ledger $\Ledger$ under
  \citep[Exogenous~Verification]{lasser2026exo}'s
  substrate-exclusivity discipline.
\end{enumerate}

Invariant $I_8$ (Paper~10 §3) extends this construction to
environment-side observables.  The mirror property is:
\[
  I_8: \quad \forall v \in V_{\mathrm{env}}:\
  \exists\ \Witness_v \text{ on } s_{\mathrm{env}} \neq
  s(\mathrm{agent}).
\]
The witness substrate $s_{\mathrm{env}}$ is environment-side; the
witnessed quantity $v$ is an exogenous state variable; the
observed-substrate constraint becomes $s(\mathrm{agent})$.

\textbf{Trust-model gap (Round A Q1).}  V1 said only
$s_{\mathrm{env}} \neq s(\mathrm{agent})$, which is insufficient:
the witness substrate must also be outside the adversary's
write-access / failure domain.  ``Distinct partition label'' does
not imply ``causally isolated from adversarial corruption''.  V2
adopts the stronger requirement, named in Assumption (C12.ENV-WIT)
below.

\begin{assumption}[(C12.ENV-WIT) Environment-side witness instantiation]
\label{ass:C12-EW}
The deployment satisfies the following four sub-clauses for each
$v \in \Vem$:

\textbf{(C12.PUB) Published fixed $\Vem$.}  $\Vem$ is enumerated
in deployment documentation before deployment activation; the
enumeration is fixed for the deployment epoch.  Operators may
re-enumerate across epochs (with corresponding re-audit), but not
during a single epoch.

\textbf{(C12.PART) Source/witness substrate partition.}  Each
$v \in \Vem$ has a designated witness substrate $s_{\mathrm{env}}(v)$
with $s_{\mathrm{env}}(v) \neq s(\mathrm{agent})$ and
$s_{\mathrm{env}}(v) \neq s_{\mathrm{source}}(v)$ (where
$s_{\mathrm{source}}(v)$ is the substrate generating $v$'s value).

\textbf{(C12.TRUST-WRITE) Witness substrate outside adversary
write-access domain.}  The adversary cannot tamper with $v$'s
recorded value through $s_{\mathrm{env}}(v)$.  The witness
substrate is operationally distinct from any substrate the
adversary holds write privileges on.

\textbf{(C12.TRUST-CORR) Witness substrate outside adversary
failure-correlation domain.}  Failure of $s_{\mathrm{env}}(v)$ is
not correlated with adversarial action on the agent: a successful
attack on the agent does not also corrupt the environment-side
witness via shared infrastructure.  This is operationally distinct
from (C12.TRUST-WRITE) — write-access exclusion handles
intentional tampering; failure-correlation exclusion handles
shared-infrastructure failures.

\textbf{(C12.SETUP) Trusted-setup honesty.}  The trusted-setup
ceremony for environment-side cryptographic primitives (Pedersen
commitment generators or MPC equivalent) has at least one honest
participant, mirroring Paper~5's at-least-one-honest assumption
(\citep[Exogenous~Verification]{lasser2026exo} commitment.tex).
Setup parameters are published before deployment activation.

\textbf{(C12.CAL) Per-variable calibration.}  Every $v \in \Vem$
(canonical or deployment-added) is calibrated with: (a) baseline
distribution $p_{0, v}$ (Bernoulli/Poisson/multinomial) with
support convention satisfied ($p_0 \in (\eenv, 1 - \eenv')$ open
for Bernoulli, $\lambda_0 \geq \eenv > 0$ for Poisson, or
(C5.HOEFF) clipping fallback); (b) feasible adversarial alternative
$p_{1, v}$ with threshold $\eta_v > 0$; (c) recorded-stream map per
the DPI precondition; (d) strictly positive raw or post-clipping
drift floor $\delta_v^{\mathrm{env}} > 0$.  Without (C12.CAL),
the strict-positivity of $\deltaadv^{\mathrm{env}} = \min_v
\delta_v^{\mathrm{env}}$ holds only for the canonical four
variables; (C12.CAL) extends it to deployment-added variables.
\end{assumption}

\textbf{Why (C12.ENV-WIT) is needed.}  Paper~5's trust model
(at-least-one-honest setup, witness outside adversarial substrate)
does not transfer automatically because the environment-side
substrate partition is operationally distinct (separate trusted-
setup ceremony, separate witness substrate identification, separate
ledger commitments).  Without (C12.ENV-WIT), an adversary that
controls $s_{\mathrm{env}}(v)$ could fabricate witness records;
without (C12.SETUP), the cryptographic commitments themselves can
be forged.  Both close gaps that the inline sketch glossed over.

\subsection{The detection class $\Aadv^{\mathrm{env}}$}

\begin{definition}[Environment-side adversarial class
$\Aadv^{\mathrm{env}}$]
A regime-(iii) strategy $q$ achieving its effect via environment
manipulation belongs to $\Aadv^{\mathrm{env}}$ if and only if $q$
produces an observable shift in at least one $v \in V_{\mathrm{env}}$
under the SPRT alternative versus baseline.
\end{definition}

\textbf{The named residual (R2), refined.}  R2 covers
manipulations satisfying \emph{both} conditions:
(i)~target only exogenous variables outside $\Vem$, \emph{and}
(ii)~produce no threshold-exceeding shift in any $v \in \Vem$
(direct or via causal propagation).
A manipulation of an unmonitored variable that causally shifts a
monitored $v \in \Vem$ above the SPRT detectability threshold
\emph{is} detected via that monitored projection — it is not R2.
Hence R2 is precisely the unmonitored-and-undetected class, not
``unmonitored alone''.

The residual is real: $\Vem$ is finite by deployment policy, and
adversarial environment manipulation that satisfies both
conditions above escapes detection.

\section{Per-variable KL floor derivations}

\textbf{$\Vem$ is deployment-policy-derived (Round A Q3).}  $\Vem$
is enumerated by deployment policy per (C12.PUB), with the four
canonical variables below as the illustrative \emph{minimum}
(every deployment must monitor at least these; deployments may add
more).  This v2 retires v1's implicit ``canonical four-variable''
framing.

\textbf{Reconciliation with $I_8$ enumeration.}  $I_8$ in
\texttt{invariants.tex} currently lists $\rext$, $\Delta r_K$,
$\lambda$, substrate-distinctness; deployment-tooling §8.1 lists
$\rext$, distinctness, $\lambda$, trusted-setup.  V2 reconciles by
\emph{dropping} $\Delta r_K$ from $\Vem$ ($\Delta r_K$ is
strategy-dependent — agent-side observable, not exogenous; it
belongs in agent-side $\Aadv$ machinery via the four-channel
agent-side ledger of Lemma~5b) and \emph{adding} trusted-setup
status flag (matching deployment-tooling §8.1).  The integration
edits $I_8$'s enumeration accordingly.

The illustrative minimum-canonical $\Vem$:
\begin{itemize}
\item $V_1$: cross-substrate cooperative novelty rate $\rext$
  (Poisson).
\item $V_2$: substrate-distinctness independence indicator
  (Bernoulli).
\item $V_3$: adversarial-event arrival rate $\lambda$ (Poisson).
\item $V_4$: trusted-setup status flag (Bernoulli).
\end{itemize}

\textbf{Support conventions (Round A Q2).}
\begin{itemize}
\item \textbf{Bernoulli} (variables $V_2, V_4$): baseline
  $p_0 \in (0, 1)$ open (excluding endpoints), with deployment-class
  lower bound $p_0 \geq \eenv > 0$ and upper bound $p_0 \leq 1 -
  \eenv$ enforced by the deployment's calibration policy.  The
  Bernoulli KL is then well-defined and bounded.
\item \textbf{Poisson} (variables $V_1, V_3$): baseline
  $\lambda_0 \geq \eenv > 0$, with $\eenv$ a deployment-class
  regularization constant.  Edge case $\lambda_0 = 0$ (no baseline
  arrivals) is handled by the regularization: the deployment uses
  $\max(\lambda_0^{\mathrm{measured}}, \eenv)$ as the operational
  baseline.
\item \textbf{Non-overlapping support / clipping.}  When raw KL
  diverges (e.g., $V_3$ with $\lambda_0 = 0$ and adversarial
  $\lambda_1 > 0$), the operational SPRT uses the clipped LLR per
  (C5.HOEFF), inherited via Lemma~6's drift-floor convention:
  $\deltastar$ in Lemma~6 is the post-clipping drift floor, which
  remains finite even when raw KL is unbounded.
\end{itemize}

\textbf{Marginal-to-joint KL via data processing (Round A Q2,
Round B precondition).}

\textbf{DPI precondition.}  For every $v \in \Vem$, the witness
protocol defines a fixed pre-deployment measurable recorded stream:
$v$'s witness $\Witness_v$ records to a fixed coordinate (or fixed
Markov kernel from joint state to recorded stream) that is
specified before deployment activation per (C12.PUB) and does not
change during the deployment epoch.  The per-variable KL
$D_{\mathrm{KL}}(q_{v_i} \,\|\, p_{0, v_i})$ is computed on this
recorded stream.  Canonical variables ($V_1, V_2, V_3, V_4$) satisfy
this precondition automatically; deployment-specific additions to
$\Vem$ must verify it (Audit~7 / Audit~8 calibration).

\textbf{DPI step.}  For any $q \in \Aadv^{\mathrm{env}}$ producing a
threshold-exceeding shift in some $v_i \in \Vem$, the marginal
distribution $q_{v_i}$ on the witness's recorded $v_i$-stream
satisfies $D_{\mathrm{KL}}(q_{v_i} \,\|\, p_{0, v_i}) \geq
\delta_i^{\mathrm{env}}$ (as derived per-variable below).  By the
data-processing inequality (DPI) applied to the marginalization
map / Markov kernel $q \mapsto q_{v_i}$:
\[
  D_{\mathrm{KL}}(q \,\|\, p_0^{\mathrm{env}}) \;\geq\;
  D_{\mathrm{KL}}(q_{v_i} \,\|\, p_{0, v_i}) \;\geq\;
  \delta_i^{\mathrm{env}}.
\]
Therefore
$D_{\mathrm{KL}}(q \,\|\, p_0^{\mathrm{env}}) \geq
\deltaadv^{\mathrm{env}} := \min_{v \in \Vem}
\delta_v^{\mathrm{env}}$.

For each variable below, we (a) state the baseline distribution
$p_{0, v}$, (b) the least-favorable adversarial alternative $p_{1, v}$
subject to the SPRT-detectability constraint, and (c) the resulting
KL divergence floor (one-sided form; two-sided variant noted at end
of §3).

\subsection{Variable $V_1$: cross-substrate cooperative rate
(Poisson)}

\textbf{Baseline.}  Under stance S0,
\citep[Horizon~Aware]{lasser2026horizon}'s cross-substrate
cooperative novelty rate $\rext > 0$ is environment-observable:
ledger entries record cross-substrate cooperative events with
Poisson arrival rate $\lambda_0^{(1)}$ in the no-adversary baseline.
The witness $\Witness_{V_1}$ on $s_{\mathrm{env}}$ records each
arrival with timestamp and substrate-pair attestation.

\textbf{Adversarial alternative.}  An environment-manipulating
strategy that suppresses cross-substrate cooperatives (e.g., by
disrupting environmental conditions for cross-substrate cooperation)
reduces the rate to at most $\lambda_0^{(1)} - \eta_1$ for some
$\eta_1 \in (0, \lambda_0^{(1)})$, where $\eta_1$ is the
SPRT-detectability threshold.

\textbf{KL divergence.}  The Poisson KL between rate-$\lambda_1$ and
rate-$\lambda_0$ processes per unit time is:
\[
  D_{\mathrm{KL}}(\mathrm{Poi}(\lambda_1) \,\|\,
  \mathrm{Poi}(\lambda_0)) \;=\; \lambda_1 \log(\lambda_1 /
  \lambda_0) - \lambda_1 + \lambda_0.
\]
This is convex in $\lambda_1$ and increasing in
$|\lambda_1 - \lambda_0|$.  The least-favorable choice within the
constraint $\lambda_1 \leq \lambda_0^{(1)} - \eta_1$ is
$\lambda_1 = \lambda_0^{(1)} - \eta_1$, giving:
\[
  \delta_1^{\mathrm{env}} \;=\; (\lambda_0^{(1)} - \eta_1)
  \log\!\frac{\lambda_0^{(1)} - \eta_1}{\lambda_0^{(1)}} -
  (\lambda_0^{(1)} - \eta_1) + \lambda_0^{(1)} \;>\; 0.
\]

\subsection{Variable $V_2$: substrate-distinctness independence
indicator (Bernoulli)}

\textbf{Baseline.}  Per-period substrate-distinctness independence
indicator $\xi_t \in \{0, 1\}$ is recorded by $\Witness_{V_2}$:
$\xi_t = 1$ if the period-$t$ failure-correlation audit verifies
$\mindep \geq m^*$ (independence holds), $\xi_t = 0$ otherwise.
Baseline rate: $\Pr[\xi_t = 1] = p_0^{(2)} \in (\eenv, 1 - \eenv')
\subset (0, 1)$ open (typically near $1$ for well-anchored
deployments; $\eenv, \eenv' > 0$ enforce well-definedness of the
Bernoulli KL via the support convention in §3).

\textbf{Adversarial alternative.}  An environment-manipulating
strategy that erodes substrate distinctness (e.g., by introducing
correlated failure modes across nominally distinct substrates)
reduces the success rate to at most $p_0^{(2)} - \eta_2$ for some
$\eta_2 \in (0, p_0^{(2)}]$.

\textbf{KL divergence.}  Bernoulli KL:
\[
  \delta_2^{\mathrm{env}} \;=\; (p_0^{(2)} - \eta_2)
  \log\!\frac{p_0^{(2)} - \eta_2}{p_0^{(2)}} +
  (1 - p_0^{(2)} + \eta_2) \log\!\frac{1 - p_0^{(2)} + \eta_2}
  {1 - p_0^{(2)}} \;>\; 0.
\]
This is the least-favorable choice within the constraint
$p_1 \leq p_0^{(2)} - \eta_2$.

\subsection{Variable $V_3$: adversarial-event arrival rate
(Poisson)}

\textbf{Baseline.}  Adversarial events arrive on the environment
side with Poisson rate $\lambda_0^{(3)} \geq 0$ (the baseline
threat-model arrival rate, calibrated by the deployment's threat
model).  Witness $\Witness_{V_3}$ records arrivals.

\textbf{Adversarial alternative.}  An adversarial strategy that
\emph{increases} the arrival rate to $\lambda_0^{(3)} + \eta_3$
(introducing additional adversarial events) is detectable if
$\eta_3 > 0$ exceeds the SPRT detection threshold.

\textbf{KL divergence.}  Same Poisson form as $V_1$, but with
positive shift $\eta_3$:
\[
  \delta_3^{\mathrm{env}} \;=\; (\lambda_0^{(3)} + \eta_3)
  \log\!\frac{\lambda_0^{(3)} + \eta_3}{\lambda_0^{(3)}} -
  (\lambda_0^{(3)} + \eta_3) + \lambda_0^{(3)} \;>\; 0.
\]
\textbf{Edge case $\lambda_0^{(3)} = 0$ (Round A Q2 fix).}  When
the deployment claims zero baseline adversarial-event arrivals
(threat-model-pristine baseline), the raw Poisson KL formally
diverges: $\eta_3 \log(\eta_3/0) = +\infty$.  Two operationally
sound resolutions:

\emph{(a) Regularized baseline (preferred):}  Adopt the support
convention from §3 and use $\lambda_0^{(3)} \geq \eenv > 0$, with
$\eenv$ a deployment-class regularization constant.  The KL becomes
finite and strictly positive.  The $\eenv$ inflation is conservative
(slightly weaker detection guarantee), in exchange for a
well-defined floor.

\emph{(b) Clipped LLR (fallback for hard $\lambda_0 = 0$):}  Use
the operational SPRT's clipped LLR per (C5.HOEFF).  Lemma~6's
post-clipping drift floor $\deltastar$ remains finite and bounded
below by the clip-radius constant $\Bclip$, independently of
whether raw KL diverges.  This handles non-overlapping support
without requiring baseline regularization.

Either resolution is acceptable; the deployment chooses based on
operational data.  The integration must surface $\eenv$ as a
deployment-class constant (added under (C5.SUPP) extended for
environment-side observables, or under (C12.ENV-WIT) sub-clause).

\subsection{Variable $V_4$: trusted-setup status flag (Bernoulli)}

\textbf{Baseline.}  The trusted-setup status flag $\sigma_t \in
\{0, 1\}$ records whether the cryptographic trusted-setup ceremony
remains intact at period $t$: $\sigma_t = 1$ if integrity audits
pass, $\sigma_t = 0$ if compromise is detected.  Baseline rate
$p_0^{(4)} \in (1 - \eenv, 1 - \eenv') \subset (0, 1)$ open
(near $1$ for properly maintained setups; $\eenv, \eenv' > 0$
deployment-class regularization constants ensuring the Bernoulli
KL is well-defined).

\textbf{Adversarial alternative.}  Environment-side adversarial
strategies that compromise the trusted-setup substrate reduce the
flag's success rate to $p_0^{(4)} - \eta_4$.

\textbf{KL divergence.}  Same Bernoulli form as $V_2$:
\[
  \delta_4^{\mathrm{env}} \;=\; (p_0^{(4)} - \eta_4)
  \log\!\frac{p_0^{(4)} - \eta_4}{p_0^{(4)}} +
  (1 - p_0^{(4)} + \eta_4) \log\!\frac{1 - p_0^{(4)} + \eta_4}
  {1 - p_0^{(4)}} \;>\; 0.
\]

\subsection{Two-sided floor option (Round A Q2)}

The per-variable derivations above use \emph{one-sided} shifts:
$V_1$ and $V_2$ degrade (downward shift in cooperative/distinctness
rates), $V_3$ inflates (upward shift in adversarial arrivals), $V_4$
degrades (downward shift in trusted-setup integrity).  These
directions match the adversarial model: degradation of audit-pass
indicators, inflation of attack rates.

If a deployment defines $\Aadv^{\mathrm{env}}$ to include
\emph{any} threshold-exceeding observable shift (rather than
direction-specific degradation), then $\delta_v^{\mathrm{env}}$
becomes the two-sided minimum:
\[
  (\delta_v^{\mathrm{env}})_{\mathrm{2s}} \;:=\; \min\bigl\{
  D_{\mathrm{KL}}(p_{0,v} - \eta_v \,\|\, p_{0,v}),\;
  D_{\mathrm{KL}}(p_{0,v} + \eta_v \,\|\, p_{0,v})
  \bigr\}.
\]
For Bernoulli with $p_0$ away from $0.5$, the directions are
asymmetric (smaller KL toward $0.5$); the minimum is the
direction-toward-$0.5$ shift.  For Poisson, larger
$|\lambda_1 - \lambda_0|$ gives larger KL, but the convexity
asymmetry around $\lambda_0$ matters: shifting up vs.\ down by the
same $\eta$ gives different KL values.

The two-sided variant is conservative (stricter for adversaries who
can choose direction) but tighter (lower $\delta^{\mathrm{env}}$
floor, hence larger $\Tbeta$ in Lemma~6).  Deployments choose based
on threat model: degradation-only or arbitrary-direction.

\subsection{Combined floor}

\[
  \deltaadv^{\mathrm{env}} \;:=\; \min_{v \in \Vem}
  \delta_v^{\mathrm{env}} \;>\; 0,
\]
where $\delta_v^{\mathrm{env}}$ is either the one-sided or
two-sided form per the deployment's threat model.  Each
$\delta_v^{\mathrm{env}} > 0$ under the support conventions of §3
(open Bernoulli baselines, regularized Poisson baselines, or
clipping for hard non-overlap); the minimum over a finite-
cardinality $\Vem$ (per (C12.PUB)) is therefore strictly positive.

\section{Statement of Lemma 5e (full)}

\begin{lemma}[Environment-side witness extension]
\label{lem:5e}
Under invariant $I_8$ (environment-side substrate-exclusive
witnesses on exogenous variables), Lemma~5b's channel-restricted
detection class extends to environment-side observables.
Specifically, for any $q \in \Aadv^{\mathrm{env}}$:
\[
  D_{\mathrm{KL}}(q \,\|\, p_0^{\mathrm{env}}) \;\geq\;
  \deltaadv^{\mathrm{env}} \;>\; 0,
\]
where $\deltaadv^{\mathrm{env}} = \min\{\delta_v^{\mathrm{env}} :
v \in V_{\mathrm{env}}\}$ is the minimum over per-variable KL
floors derived from environment-side least-favorable alternatives.

For the canonical four-variable $V_{\mathrm{env}} = \{V_1, V_2,
V_3, V_4\}$ enumerated in §3, the explicit per-variable floors are:
\begin{itemize}
\item $\delta_1^{\mathrm{env}}$ (Poisson cross-substrate cooperative
  rate),
\item $\delta_2^{\mathrm{env}}$ (Bernoulli substrate-distinctness
  indicator),
\item $\delta_3^{\mathrm{env}}$ (Poisson adversarial-event arrival
  rate),
\item $\delta_4^{\mathrm{env}}$ (Bernoulli trusted-setup status
  flag),
\end{itemize}
each $> 0$ as derived in §3.1--3.4.
\end{lemma}

\begin{proof}
Four parts: (a) the symmetric construction with trust-model
verification (§2 + (C12.ENV-WIT)), (b) the per-variable KL floor
derivations under the support conventions (§3.1--3.4), (c) the
DPI-based combined-floor argument, (d) the strict-positivity of
the minimum.

\emph{(a) Symmetric construction.}  Under (C12.ENV-WIT), each
$v \in \Vem$ has a witness substrate $s_{\mathrm{env}}(v)$ with
the four sub-clauses (PUB published $\Vem$, PART partition,
TRUST witness-outside-adversary, SETUP honest setup).  Paper~5's
cryptographic and ledger machinery transfers under (C12.SETUP);
the trust model is closed by (C12.TRUST).

\emph{(b) Per-variable derivations.}  For each $v_i \in \Vem$,
under the support conventions ($p_0 \in (\eenv, 1 - \eenv')$ for
Bernoulli, $\lambda_0 \geq \eenv > 0$ for Poisson, or (C5.HOEFF)
clipping for non-overlapping support), the least-favorable
adversarial alternative within the SPRT-detectability constraint
gives a strictly positive per-variable KL floor
$\delta_i^{\mathrm{env}} > 0$ (§3.1--3.4 derivations; one-sided or
two-sided per the deployment's threat model).

\emph{(c) Combined-floor via DPI.}  For any $q \in
\Aadv^{\mathrm{env}}$ producing a threshold-exceeding shift in some
$v_i \in \Vem$, the marginal-to-joint KL step (§3) gives
$D_{\mathrm{KL}}(q \,\|\, p_0^{\mathrm{env}}) \geq
D_{\mathrm{KL}}(q_{v_i} \,\|\, p_{0, v_i})$ via the data-processing
inequality applied to the marginalization map.  Combined with the
per-variable floor $\delta_i^{\mathrm{env}}$:
\[
  D_{\mathrm{KL}}(q \,\|\, p_0^{\mathrm{env}}) \;\geq\;
  \delta_i^{\mathrm{env}} \;\geq\; \deltaadv^{\mathrm{env}} :=
  \min_{v \in \Vem} \delta_v^{\mathrm{env}}.
\]

\emph{(d) Strict positivity.}  $\Vem$ is finite by (C12.PUB);
each $\delta_v^{\mathrm{env}} > 0$ by part (b); therefore
$\deltaadv^{\mathrm{env}} > 0$.  $\Box$
\end{proof}

\section{Discussion}

\paragraph{Composition with Lemma 6.}  The post-clipping union-class
drift floor $\deltastar$ used in Lemma~6's $\Tbeta$ derivation is
$\deltastar \leq \min(\deltaadv, \deltaadv^{\mathrm{env}})$, with
the inequality accounting for any drift reduction from (C5.HOEFF)
clipping.  Lemma~5e supplies $\deltaadv^{\mathrm{env}}$ for the
union; Lemma~5b supplies $\deltaadv$.

\paragraph{Coverage and the named residual (R2), refined form.}
$\Vem$ is finite by deployment policy (C12.PUB).  Variables outside
$\Vem$ are real exogenous quantities not monitored by
environment-side witnesses.  R2 covers the precise class of
manipulations satisfying \emph{both}:
(i)~target only exogenous variables outside $\Vem$, \emph{and}
(ii)~produce no threshold-exceeding shift in any monitored
$v \in \Vem$ (direct or via causal propagation).  Manipulations of
unmonitored variables that causally shift a monitored $v$ are
detected through the monitored projection — \emph{not} R2.

The deployment must enumerate $\Vem$ to cover its claimed material
environment threat model; coverage gaps shift adversarial mass into
the precise R2 class (the and-condition above), not into
``everything outside $\Vem$''.

\paragraph{Asymmetric vs. symmetric machinery.}
Paper~5's agent-side substrate-exclusivity is constructed; Paper~10
§3 specifies that environment-side substrate-exclusivity is the
mirror, but is new machinery (separate substrate partition, separate
trust models).  This lemma's symmetric-construction argument is
\emph{structural}: it inherits Paper~5's discipline.  The
deployment must instantiate the environment-side construction
(Audit~2/3 in §8 deployment tooling).

\paragraph{Choice of $\eta_v$.}
Each per-variable shift threshold $\eta_v$ is the SPRT
detectability threshold for that variable.  Operationally, $\eta_v$
is set by the deployment's threat model: how much shift in $v$
must be detected for the deployment to consider the manipulation
adversarially material.  Tighter $\eta_v$ gives smaller
$\delta_v^{\mathrm{env}}$ (higher detection sensitivity but lower
KL floor); looser $\eta_v$ gives larger $\delta_v^{\mathrm{env}}$
(lower sensitivity but higher floor, hence faster detection in
$\Tbeta$).

\section{What this proof does and does not establish}

\textbf{Establishes:}
\begin{itemize}
\item The symmetric-construction argument formalizing $I_8$'s
  extension of Paper~5's substrate-exclusivity to environment-side
  observables.
\item Per-variable KL floor derivations for the four canonical
  environment-side observables (Poisson for rates; Bernoulli for
  binary flags), each strictly positive.
\item The combined floor $\deltaadv^{\mathrm{env}} = \min_v
  \delta_v^{\mathrm{env}} > 0$ as a deployment-class constant.
\item Composition with Lemma~6 via the post-clipping union-class
  drift floor $\deltastar$.
\end{itemize}

\textbf{Does not establish (deferred):}
\begin{itemize}
\item Tightness of per-variable shifts $\eta_v$.  These are
  threat-model-derived and deployment-policy-set.
\item Exhaustive coverage of $V_{\mathrm{env}}$.  This is a
  deployment-policy enumeration question; under-coverage falls into
  residual~(R2).
\item Security proof of the environment-side substrate-exclusivity
  construction.  Paper~10 §3 specifies the construction, not the
  security proof; cryptographic security is inherited from
  Paper~5's substrate-exclusivity discipline.
\item Multinomial extension.  $V_{\mathrm{env}}$ is enumerated as
  Poisson + Bernoulli observables for the canonical four; if
  deployment policy adds multinomial environment-side observables
  (e.g., partition over trade-flow categories), the analogous
  multinomial KL derivation from Lemma~5b's appendix proof applies.
\end{itemize}

\section{Specific verification questions for v2 codex review}

\begin{enumerate}
\item \textbf{(Q1 trust model.)}  V2 introduces (C12.ENV-WIT) with
  four sub-clauses: (PUB) published fixed $\Vem$, (PART) source/
  witness substrate partition, (TRUST) witness substrate outside
  adversary's write-access/failure domain, (SETUP) trusted-setup
  honesty.  Is this the right factoring?  In particular: should
  (TRUST) be split further (write-access vs failure-correlation
  are distinct properties), or is the bundling cleaner?  Does
  (SETUP) need to repeat Paper~5's at-least-one-honest assumption
  in full, or is the cross-reference sufficient?

\item \textbf{(Q2 KL derivation.)}  V2 adds the marginal-to-joint
  DPI step + support conventions ($p_0 \in (\eenv, 1-\eenv')$
  Bernoulli; $\lambda_0 \geq \eenv$ Poisson; clipping fallback) +
  two-sided floor option.  Is the DPI argument tight?  In
  particular: the DPI requires that the witness-recorded
  observable $v_i$ is a deterministic function of the joint
  state, which holds for the canonical observables (rates +
  Bernoulli flags) but should be verified for any deployment-
  specific additions.  Does the proof need to make this DPI
  precondition explicit?

\item \textbf{(Q3 $\Vem$ enumeration.)}  V2 makes $\Vem$
  deployment-policy-derived with the four canonical variables as
  illustrative minimum, and reconciles the $I_8$ vs deployment-
  tooling mismatch by dropping $\Delta r_K$ from $\Vem$ and adding
  trusted-setup status.  Is this reconciliation right?  Specifically:
  is $\Delta r_K$ correctly classified as agent-side (not
  environment-side), and does removing it from $I_8$'s enumeration
  cause issues elsewhere in the paper (e.g., does any other proof
  rely on $\Delta r_K$ being witnessed by $I_8$)?

\item \textbf{(Q4 R2 refinement.)}  V2 refines R2: ``manipulations
  targeting only unmonitored exogenous variables \emph{and}
  producing no threshold-exceeding monitored shift.''  Is this
  refinement consistent with Theorem~1's R2 statement?  Should
  the integration update Theorem~1's R2 wording to match this
  precise form?

\item \textbf{(Q5 theorem-level conditions.)}  V2 introduces
  (C12.ENV-WIT) and notes that (C5.SPRT) sub-clauses should
  explicitly cover environment-side LLR.  Is the new (C12) the
  right framing, or should the conditions live as expanded $I_8$
  sub-clauses?  Pattern from prior lemmas (C9 BD/CL/WB/TF, C10
  CN/SU, C5.SPRT HOEFF/MULT/IID/SUPP) supports a new
  (C12.ENV-WIT) bundle; is the bundle's sub-clause structure
  parallel enough?

  Additionally: V2 adds $\eenv$ as a deployment-class regularization
  constant for support conventions.  Should $\eenv$ surface as a
  separate (C5.SUPP) sub-clause, or is folding it into Audit~7
  (which already covers (C5.SPRT) calibration) sufficient?
\end{enumerate}

\end{document}
